% SPDX-License-Identifier: LPPL-1.3c
% Copyright (C) 2026 Christophe Jorssen
% Author and Current Maintainer: Christophe Jorssen <christophe.jorssen@gmail.com>
% LPPL maintenance status: maintained
% This file is part of luafemm. See LICENSE and LICENSES.md for its terms.
\section{Electrical networks and IEC schematics}
\label{sec:electrical-networks}

A winding is both a magnetic source and an electrical port. Connecting a resistor
in series with it changes its current; placing a second winding nearby changes
its harmonic terminal voltage even without an electrical connection. A network
therefore shares its winding ports with the field calculation. Its wires describe
logical connections, independently of the coordinates used to draw them.

The electrical model is available through the main |femm| library. Load the
optional |femm.circuits.ee.IEC| library to declare data on IEC symbols or display
those data in an electrical schematic. The interface uses native TikZ |node|,
|to|, |info| and symbol styles. An ordinary unannotated symbol has no physical
effect. The following loading instruction is also executed for this chapter.
\begin{codeexample}[setup code]
\usetikzlibrary{femm.circuits.ee.IEC}
\end{codeexample}

\subsection{Connect a source, a resistance and a winding}
Declare the circuit, then the network, then its elements. The terminals |p| and
|n| are electrical identifiers: the voltage is $U=V(p)-V(n)$ and positive current
flows from |p| to |n|. Repeating an identifier makes a connection; crossing two
lines in a drawing does not. Node |0| is the default voltage reference.

\luafemmexamplebegin{electrical-network}
\begin{codeexample}[width=0pt]
\begin{tikzpicture}[
  femm/problem={name=network-demo,frequency=200,
    xmin=-10,xmax=10,ymin=-8,ymax=8,mesh size=1,depth=20},
  femm/circuit={name=coil,drive=network},
  femm/network={name=supply},
  femm/element={name=source,kind=voltage source,p=a,n=0,value=1},
  femm/element={name=load,kind=resistor,p=a,n=b,value=2},
  femm/element={name=winding,kind=winding,p=b,n=0,circuit=coil},
  scale=2]
  \filldraw[fill=orange!25,
    femm/region={material=copper,circuit=coil,turns=20}]
    (-6,-4) rectangle (-4,4);
  \filldraw[fill=orange!25,
    femm/region={material=copper,circuit=coil,turns=-20}]
    (4,-4) rectangle (6,4);
  \femmmeasure[name=coil-terminal,quantity=circuit,circuit=coil]
  \femmmeasure[name=coil-element,quantity=network,element=winding]
  \pic[femm/lines=10] {femm field};
\end{tikzpicture}
\end{codeexample}

The winding symbol below refers to the already computed field port. It does not
add a lumped inductance or count copper resistance a second time. The external
$2\,\Omega$ resistor is a distinct component. The source's prescribed voltage
is a peak amplitude; its phase defaults to zero.
\begin{codeexample}[width=0pt]
\begin{tikzpicture}[circuit ee IEC]
  \draw (0,0) to[voltage source={info={$1\,\mathrm V$}}] (0,2)
    to[resistor={info={$2\,\Omega$}}] (4,2)
    to[inductor={femm/ee={model=network-demo,
      element=winding,result=coil-element,values=solution,
      show=current,representation=rms}}] (4,0) -- (0,0);
\end{tikzpicture}
\end{codeexample}

A canonical branch view is useful when the topology was declared without a
schematic. Repeated terminal labels denote the same electrical node. The branch
view keeps each connection explicit without an automatic wire-routing heuristic.
Use ordinary TikZ placement when a conventional loop or bridge drawing is wanted.
\begin{codeexample}[]
\begin{tikzpicture}
  \pic[femm/ee view={model=network-demo}] {femm ee circuit};
\end{tikzpicture}
\end{codeexample}
\luafemmexampleend{electrical-network}

\subsection{Let a symbol declare the excitation}
Declarations are optional. Set |mode=declare| on a symbol to forward its physical
options to the same circuit or element interface used above. Only the symbol
that owns the declaration uses this mode; other views use |mode=view|. Repeating
a declaration is an error, rather than an implicit change of the source.

Finish all physical declarations before a field picture, measurement or solve.
Afterwards, results may be displayed in several pictures. Rendering a view does
not start a calculation. This separation avoids solving a half-declared network
when TikZ typesets an early label.
\luafemmexamplebegin{iec-declaration}
\begin{codeexample}[width=0pt]
\begin{tikzpicture}[circuit ee IEC,
  femm/problem={name=iec-coil,xmin=-10,xmax=10,
    ymin=-8,ymax=8,mesh size=1,depth=20}]
  \draw (35,-8) to[current source={femm/ee={mode=declare,
    circuit=drive,circuit options={current=2},show=current}}] (35,8);
  \path[femm/region={material=copper,circuit=drive,turns=20}]
    (-6,-4) rectangle (-4,4);
  \path[femm/region={material=copper,circuit=drive,turns=-20}]
    (4,-4) rectangle (6,4);
  \femmmeasure[name=iec-result,quantity=circuit,circuit=drive]
  \pic[femm/lines=8] {femm field};
\end{tikzpicture}
\end{codeexample}

The reverse direction uses the same circuit declaration and a saved measurement.
Its source and winding are generated as a simple loop. All distances here are
schematic distances; they do not change the winding geometry.
\begin{codeexample}[]
\begin{tikzpicture}
  \pic[femm/ee view={model=iec-coil,circuit=drive,
    result=iec-result,values=solution,show=voltage}] {femm ee circuit};
\end{tikzpicture}
\end{codeexample}
\luafemmexampleend{iec-declaration}

\subsection{Prescribe flux linkage and determine current}
A current drive specifies $I$ and computes $\lambda$. A linkage drive specifies
$\lambda$ and solves for $I$ together with the field. This inverse problem retains
nonlinear magnetic materials and permanent-magnet offsets. It does not divide
the target by a constant inductance. The displayed current below is computed.
\luafemmexamplebegin{prescribed-linkage}
\begin{codeexample}[width=0pt]
\begin{tikzpicture}[
  femm/problem={xmin=-10,xmax=10,ymin=-8,ymax=8,
    mesh size=1,depth=20},
  femm/circuit={name=controlled,drive=linkage,linkage=2e-5},scale=2]
  \filldraw[fill=gray!25,femm/region={material=armco-pure-iron}]
    (-2,-5) rectangle (2,5);
  \filldraw[fill=orange!25,
    femm/region={material=copper,circuit=controlled,turns=20}]
    (-6,-4) rectangle (-4,4);
  \filldraw[fill=orange!25,
    femm/region={material=copper,circuit=controlled,turns=-20}]
    (4,-4) rectangle (6,4);
  \femmmeasure[name=controlled-terminal,quantity=circuit,circuit=controlled]
  \pic[femm/lines=10] {femm field};
  \node[above] at (0,9) {$I=\femmvalue[component=current]{controlled-terminal}$};
\end{tikzpicture}
\end{codeexample}
\luafemmexampleend{prescribed-linkage}

\subsection{Electrical declaration reference}
\label{sec:network-keys}
All declarations belong to the current magnetic model and precede meshing.
Identifiers contain letters, digits, underscores or hyphens. One network per
model is supported; every node must be connected topologically to its reference.
Floating electrical systems, contradictory ideal sources and singular constraint
systems produce errors. A circuit port may appear in only one winding element.

\begin{command}{\femmnetwork\oarg{options}}
Declare a named network with the keys below. The equivalent picture/path key is
|/tikz/femm/network={options}|; picture-level declarations execute in their
written order when the model starts.
\end{command}
\begin{key}{/tikz/femm/network=\marg{options}}
Equivalent to |\femmnetwork|. The key has no physical effect until applied.
\end{key}
\begin{key}{/femm/network/name=\meta{identifier} (initially empty)}
Required network name, local to the magnetic model.
\end{key}
\begin{key}{/femm/network/ground=\meta{node} (initially 0)}
Electrical reference node with zero scalar voltage. This is separate from the
magnetic potential gauge and does not impose $A_z=0$ at a geometric point.
\end{key}
\begin{command}{\femmelement\oarg{options}}
Declare one oriented electrical branch. The same options can be supplied to
|/tikz/femm/element|. Command defaults are reset on each invocation;
PGF styles may supply shared parameters.
\end{command}
\begin{key}{/tikz/femm/element=\marg{options}}
Forward an element declaration. It may be placed on an IEC symbol even without
the optional adapter; native graphical labels remain independently customisable.
\end{key}
\begin{key}{/femm/element/name=\meta{identifier} (initially empty)}
Required unique branch name within the network.
\end{key}
\begin{key}{/femm/element/kind=\meta{choice} (initially resistor)}
One of |resistor|, |inductor|, |capacitor|, |current source|, |voltage source|,
or |winding|. A winding binds a complete named FE circuit, including its supported
series/parallel internal branches. Other kinds are ideal lumped components.
\end{key}
\begin{key}{/femm/element/p=\meta{node} (initially empty)}
Required positive voltage terminal and positive-current entry terminal.
\end{key}
\begin{key}{/femm/element/n=\meta{node} (initially empty)}
Required negative voltage terminal, distinct from |p|.
\end{key}
\begin{key}{/femm/element/value=\meta{number} (initially 0)}
Resistance in ohms, inductance in henries, capacitance in farads, source current
in amperes or source voltage in volts. Passive values must be positive; source
values may be signed. At nonzero frequency, source values are peak amplitudes.
A winding requires zero: its impedance is calculated from the FE model.
\end{key}
\begin{key}{/femm/element/phase=\meta{degrees} (initially 0)}
Temporal phase of an independent harmonic source. Passive components and winding
bindings require zero. A DC problem also requires zero.
\end{key}
\begin{key}{/femm/element/circuit=\meta{identifier} (initially empty)}
Required for |kind=winding| and forbidden for other kinds. The named circuit
must use |drive=network|. Its material conductivity and conductor length must
permit a finite positive winding resistance.
\end{key}

\begin{key}{/femm/circuit/drive=\meta{choice} (initially current)}
In addition to |current| and |voltage|, use |network|, |linkage| or |flux|.
A network drive gets its current and voltage from its electrical connection.
Linkage and flux drives require a series circuit; planar windings must close
their signed turns. Prescribe only the selected excitation quantity; unused
|current|, |voltage|, |linkage| and |flux| values must be zero.
\end{key}
\begin{key}{/femm/circuit/linkage=\meta{weber} (initially 0)}
Target flux linkage $\lambda=\sum N_k\Phi_k$ for |drive=linkage|. The static value
is signed; the harmonic value uses the circuit's peak amplitude and |phase|.
\end{key}
\begin{key}{/femm/circuit/flux=\meta{weber} (initially 0)}
Target mean linked flux per reference turn for |drive=flux|. The imposed
constraint is $\lambda=N_{\rm ref}\Phi$. This is a winding observable, not a
claim that leakage vanishes or that every geometric cross section carries the
same flux. Use |flux turns| to define $N_{\rm ref}$ explicitly.
\end{key}
\begin{key}{/femm/circuit/flux turns=\meta{positive number} (initially 1)}
Reference turn count for the preceding conversion. It does not replace the
signed region |turns| used to excite the field or compute winding resistance.
\end{key}

\subsection{Measurements and plotting}
\begin{key}{/femm/measure/quantity=network}
Measure the electrical network after solving. With neither |element| nor |node|,
the result includes the signed sum of absorbed powers and absolute residuals
|kcl_residual| (A) and |branch_residual| (mixed branch-law units, reported as a
maximum diagnostic rather than a physical scalar). Field convergence is checked
separately. Use named results with |\femmvalue|, |\femmresult| and
|\femmresultplot| exactly as for circuit measurements.
\end{key}
\begin{key}{/femm/measure/element=\meta{identifier} (initially empty)}
For |quantity=network|, select branch current, voltage, absorbed |power| and
|reactive_power|. Harmonic components use the suffixes |_real|, |_imaginary|,
|_amplitude|, |_rms| and |_phase|. A zero phasor has no phase. A source supplying
energy has negative absorbed power under the passive reference convention.
\end{key}
\begin{key}{/femm/measure/node=\meta{identifier} (initially empty)}
For |quantity=network|, select node voltage relative to |ground|. Mutually
exclusive with |element|. DC branch results also provide |current| and |voltage|.
\end{key}

\subsection{IEC view reference}
\label{sec:ee-keys}
\begin{key}{/tikz/femm/ee=\marg{options} (default empty)}
Bind or declare a physical item on a native IEC symbol. The options below inherit
in the current TikZ scope. The generated label uses native |info|; a later |info|
option can replace or reposition it. Reusable |every inductor| or custom TikZ
styles may bind views, but a declaration must occur exactly once.
\end{key}
\begin{key}{/tikz/femm/ee view=\marg{options}}
Configure a passive generated pic without declaring a physical object or placing
an |info| label on the surrounding scope.
\end{key}
\begin{key}{/femm/ee/mode=\meta{choice} (initially view)}
|view| reads existing data. |declare| forwards |circuit options| or |element options| before displaying the item. An ordinary IEC symbol without |femm/ee|
never registers an electrical object.
\end{key}
\begin{key}{/femm/ee/model=\meta{name} (initially empty)}
Select a named model. If empty, the named result supplies its model; otherwise
the current model is used. An explicit model and result must agree.
\end{key}
\begin{key}{/femm/ee/circuit=\meta{name} (initially empty)}
Bind an existing winding circuit, or name a circuit to declare. Mutually exclusive
with |element|. Names in automatically generated labels are safely typeset.
\end{key}
\begin{key}{/femm/ee/element=\meta{name} (initially empty)}
Bind or declare a network element. It must belong to the selected model.
\end{key}
\begin{key}{/femm/ee/circuit options=\marg{keys} (initially empty)}
Options forwarded to |\femmcircuit| in declaration mode. The view's |circuit|
provides the name. Use the ordinary circuit key family and its styles.
\end{key}
\begin{key}{/femm/ee/element options=\marg{keys} (initially empty)}
Options forwarded to |\femmelement| in declaration mode. The view's |element|
provides the name. The network must already exist.
\end{key}
\begin{key}{/femm/ee/result=\meta{name} (initially empty)}
Saved result used for solution values. A circuit view requires a measurement of
that circuit; an element view requires a measurement of that element.
\end{key}
\begin{key}{/femm/ee/values=\meta{choice} (initially input)}
|input| displays explicitly prescribed quantities. |solution| displays a completed
named measurement. An unavailable quantity is an error; labels never guess a
voltage, a current or a phase and never implicitly solve the model.
\end{key}
\begin{key}{/femm/ee/show=\marg{list} (initially name)}
Comma-separated |name|, |current|, |voltage|, |power| or |linkage|; |none| suppresses
the generated content. Several quantities share one label. Native |info| and
|info'| can provide an alternative layout using |\femmvalue|.
\end{key}
\begin{key}{/femm/ee/representation=\meta{choice} (initially amplitude)}
Harmonic display: |amplitude|, |rms|, |real|, |imaginary| or |phase|. Powers are
already cycle averages. DC quantities remain signed real values. This choice
changes only the displayed solution, never the prescribed physical source.
\end{key}
\begin{key}{/femm/ee/reference=\meta{choice} (initially forward)}
|forward| uses the physical port reference. |reverse| reverses both displayed
current and voltage references (and linkage); signed components change sign,
phase changes by $180^\circ$, magnitudes and absorbed power remain unchanged.
The drawing's own arrow orientation is controlled by native TikZ options.
\end{key}
\begin{stylekey}{/tikz/femm ee wire}
Generated wires default to |semithick|. Append colour, line width or other ordinary
TikZ drawing options without changing the calculation.
\end{stylekey}
\begin{key}{/tikz/pics/femm ee circuit}
With |circuit|, draw a current/voltage source and winding loop. With neither
|circuit| nor |element|, draw the selected model's network as a labelled branch
list. Repeated node names mean electrical connection. The latter view displays
input names; measured branch annotations can be added on hand-placed symbols.
The internal layout uses centimetre basis vectors; surrounding transformations
and ordinary TikZ styling remain available.
\end{key}

\subsection{Equations and limits}
At angular frequency $\omega$, a resistor obeys $U=RI$, an ideal inductor
$U=\mathrm i\omega LI$ and a capacitor $I=\mathrm i\omega CU$. At zero frequency,
these become an exact short circuit and open circuit respectively. No artificial
leak resistance is inserted. The graph uses explicit branch currents and nodal
voltages, while FE winding currents and induced voltages are solved with the
magnetic response. Mutual coupling, supported solid-conductor redistribution and
wire losses are retained. Peak phasors give
$\langle P\rangle=\tfrac12\Re(U\overline I)$.

Static linkage constraints use a damped Newton solve and an explicitly eliminated
constraint system. The magnetic tangent is solved with the selected field solver;
the small circuit system uses scaled pivoting. Harmonic linkage constraints use
the complex coupled response. Static magnetic nonlinearity remains supported;
nonlinear switching devices, time-domain RLC state and general nonlinear harmonic
networks are outside this model. Electrical voltages are lumped terminal quantities,
not differences of the magnetic vector potential.

Declare networks on native magnetic problems. Network-bearing models cannot be
captured as operating points or moving templates. Those captures retain their
existing prescribed-current contract.

Network topology and excitation data enter the solution cache; diagram layout and
labels do not. Prescribed-linkage currents are saved with the field. The general
network and linkage constraints have no exact representation as FEMM's ordinary
circuit properties. Export therefore rejects them; construct a separate terminal-
equivalent problem when exchanging a particular solved operating condition.
